%%%%%%%%%%%%%%%%%%%%%%%%%%%%%%%%%%%%%%%%%%%
%
% From a template maintained at https://github.com/jamesrobertlloyd/cbl-tikz-poster
%
%%%%%%%%%%%%%%%%%%%%%%%%%%%%%%%%%%%%%%%%%%%

\documentclass[portrait,a0b,final,a4resizeable]{a0poster}
\setlength{\paperwidth}{36in} % A0 width: 46.8in
\setlength{\paperheight}{48in} % A0 width: 46.8in

\usepackage{atbegshi}% http://ctan.org/pkg/atbegshi
\AtBeginDocument{\AtBeginShipoutNext{\AtBeginShipoutDiscard}}
\usepackage{qrcode}
\usepackage{multicol}
\usepackage{enumitem}
\usepackage{mathtools}
%\usepackage{color}
%\usepackage{morefloats}
%\usepackage[pdftex]{graphicx}
%\usepackage{rotating}
\usepackage{amsmath, amsthm, amssymb, bm}
%\usepackage{array}
%\usepackage{booktabs}
\usepackage{multirow}
%\usepackage{hyperref}
\usepackage{pgf-soroban}
\usepackage{bussproofs}
\usepackage{nicematrix}
\usetikzlibrary{cd,shapes.geometric,arrows,chains,matrix,positioning,scopes,calc,trees}
\tikzstyle{mybox} = [draw=white, rectangle]
%\definecolor{darkblue}{rgb}{0,0.08,0.45}
%\definecolor{blue}{rgb}{0,0,1}
%\usepackage{dsfont}
\usepackage[margin=0.5in]{geometry}
%\usepackage{fp}

%%%%%%%%%%%%%%%%%%%%%%%%%%%%%%%%%%%%%%%%%%%
%
% myfig
%
% \myfig - replacement for \figure
% necessary, since in multicol-environment
% \figure won't work
%
%%%%%%%%%%%%%%%%%%%%%%%%%%%%%%%%%%%%%%%%%%%

\newcommand{\myfig}[3][0]{
\begin{center}
    \vspace{1.5cm}
    \includegraphics[width=#3\hsize,angle=#1]{#2}
    \nobreak\medskip
\end{center}}

%%%%%%%%%%%%%%%%%%%%%%%%%%%%%%%%%%%%%%%%%%%
%
% mycaption
%
% \mycaption - replacement for \caption
% necessary, since in multicol-environment \figure and
% therefore \caption won't work
%
%%%%%%%%%%%%%%%%%%%%%%%%%%%%%%%%%%%%%%%%%%%

%\newcounter{figure}
\setcounter{figure}{1}
\newcommand{\mycaption}[1]{
\vspace{0.5cm}
\begin{quote}
{{\sc Figure} \arabic{figure}: #1}
\end{quote}
\vspace{1cm}
\stepcounter{figure}
}

%%%%%%%%%%%%%%%%%%%%%%%%%%%%%%%%%%%%%%%%%%%
%
% Some standard colours
%
%%%%%%%%%%%%%%%%%%%%%%%%%%%%%%%%%%%%%%%%%%%

\definecolor{camlightblue}{rgb}{0.601 , 0.8, 1}
\definecolor{camdarkblue}{rgb}{0, 0.203, 0.402}
\definecolor{camred}{rgb}{1, 0.203, 0}
\definecolor{camyellow}{rgb}{1, 0.8, 0}
\definecolor{lightblue}{rgb}{0, 0, 0.80}
\definecolor{white}{rgb}{1, 1, 1}
\definecolor{whiteblue}{rgb}{0.80, 0.80, 1}

%%%%%%%%%%%%%%%%%%%%%%%%%%%%%%%%%%%%%%%%%%%
%
% Some look and feel definitions
%
%%%%%%%%%%%%%%%%%%%%%%%%%%%%%%%%%%%%%%%%%%%

\setlength{\columnsep}{0.03\textwidth}
\setlength{\columnseprule}{0.0018\textwidth}
\setlength{\parindent}{0.0cm}

%%%%%%%%%%%%%%%%%%%%%%%%%%%%%%%%%%%%%%%%%%%
%
% \mysection - replacement for \section*
%
% Puts a pretty box around some text
% TODO - any other thoughts for what this box should look like
%
%%%%%%%%%%%%%%%%%%%%%%%%%%%%%%%%%%%%%%%%%%%

\tikzstyle{mysection} = [rectangle,
draw=none,
shade,
outer color=camlightblue!30,
inner color=camlightblue!30,
text width=0.965\columnwidth,
text centered,
rounded corners=20pt,
minimum height=0.09\columnwidth]

\newcommand{\mysection}[1]
{
\begin{center}
    \begin{tikzpicture}
        \node[mysection] {\sffamily\bfseries\LARGE#1};
    \end{tikzpicture}
\end{center}
}

%%%%%%%%%%%%%%%%%%%%%%%%%%%%%%%%%%%%%%%%%%%
%
% Set the font
%
% TODO - Not sure what a canonical choice is - feel free to modify
%
%%%%%%%%%%%%%%%%%%%%%%%%%%%%%%%%%%%%%%%%%%%

\renewcommand{\familydefault}{cmss}
\sffamily

%%%%%%%%%%%%%%%%%%%%%%%%%%%%%%%%%%%%%%%%%%%%%%%%%%%%
%%%               Background                     %%%
%%%%%%%%%%%%%%%%%%%%%%%%%%%%%%%%%%%%%%%%%%%%%%%%%%%%

\newcommand{\background}[3]{
%\definecolor{cgradbegin}{#1}
%\definecolor{cgradend}{#2}
% \psframe[fillstyle=gradient,gradend=cgradend,
% gradbegin=cgradbegin,gradmidpoint=#3](0.,0.)(1.\textwidth,-1.\textheight)
}




%%%%%%%%%%%%%%%%%%%%%%%%%%%%%%%%%%%%%%%%%%%%%%%%%%%%
%%%                pcolumn                       %%%
%%%%%%%%%%%%%%%%%%%%%%%%%%%%%%%%%%%%%%%%%%%%%%%%%%%%

\newenvironment{pcolumn}[1]{
\begin{minipage}{#1\textwidth}
\begin{center}
}{
\end{center}
\end{minipage}
}



%%%%%%%%%%%%%%%%%%%%%%%%%%%%%%%%%%%%%%%%%%%%%%%%%%%%
%%%                pbox                          %%%
%%%%%%%%%%%%%%%%%%%%%%%%%%%%%%%%%%%%%%%%%%%%%%%%%%%%

\definecolor{lcolor}{rgb}{0, 0, 0.80}
\definecolor{gcolor1}{rgb}{1, 1, 1}
\definecolor{gcolor2}{rgb}{.80, .80, 1}

% \def\fc{fillcolor}
% \def\getfc #1=#2\par{\def\ffc{#1} \ifx\ffc\fc #2\fi}
% \def\getfillcolor #1,#2\par{\getfc #1\par \getfc #2\par}

%  \newcommand{\psshadowbox}[2]{%[2][magenta]{
%      \fbox{Input arg: #1}
%      \fbox{#1}
%      \fbox {\getfillcolor #1\par}
%      \def\col{\getfillcolor #1\par}

%      \let\coll=\col
%       \coll
%     \colorbox{\col}{#2}
%       \mbox
%   \coloredshadowbox{black}{\coll}{#2}
%   }

\newcommand{\pbox}[4]{
%\psshadowbox[#3]{
%\fbox{
\mbox{
\begin{minipage}[t][#2][t]{#1}
#4
\end{minipage}
}%}
}

%%%%%%%%%%%%%%%%%%%%%%%%%%%%%%%%%%%%%%%%%%%
%
% Poster environment
%
% Centres everything and can be used to define the width of the content
%
%%%%%%%%%%%%%%%%%%%%%%%%%%%%%%%%%%%%%%%%%%%

\newenvironment{poster}{
\begin{center}
\begin{minipage}[c]{\textwidth}
}{
\end{minipage}
\end{center}
}

\def\newarrow{\mbox{\begin{tikzpicture}
\useasboundingbox{(-3pt,-4.5pt) rectangle (19pt,1pt)};
\draw[->] (0,-0.07)--(17pt,-0.07);\end{tikzpicture}}}

%%%%%%%%%%%%%%%%%%%%%%%%%%%%%%%%%%%%%%%%%%%
%
% Bottom box
%
%%%%%%%%%%%%%%%%%%%%%%%%%%%%%%%%%%%%%%%%%%%

\newlength{\bottomboxheight}
\setlength{\bottomboxheight}{0.1\paperheight}

\newcommand{\bottombox}[1]{\vfill
\noindent\colorbox{white}{
\begin{minipage}[c][\bottomboxheight][c]{\textwidth}
\centering
\begin{minipage}{0.9\textwidth}
\vfill{

\fontsizesection\color{black}
#1
}

\end{minipage}
\end{minipage}

}
}

%% Bottom box logo
\newcommand{\bottomboxlogo}[2][width=\textwidth]{
\begin{minipage}[c][\bottomboxheight][c]{0.3\textwidth}
\raggedleft\includegraphics[#1]{#2}
\end{minipage}
}

\newcommand{\bottomboxlogoleft}[2][width=\textwidth]{
\begin{minipage}[l][\bottomboxheight][c]{0.3\textwidth}
\raggedleft\includegraphics[#1]{#2}
\end{minipage}
}

%%%%%%%%%%%%%%%%%%%%%%%%%%%%%%%%%%%%%%%%%%%
%
% Highlighting
%
%%%%%%%%%%%%%%%%%%%%%%%%%%%%%%%%%%%%%%%%%%%

\definecolor{slightgray}{rgb}{0.90, 0.90, 0.90}

\usepackage{soul}
\makeatletter
\def\SOUL@hlpreamble{%
\setul{}{3.0ex}%
\let\SOUL@stcolor\SOUL@hlcolor%
\SOUL@stpreamble%
}
\makeatother

\newcommand{\inline}[1]{%
\begingroup%
\sethlcolor{slightgray}%
\hl{\ttfamily\small #1}%
\endgroup
}

\newcommand{\tinline}[1]{%
\begingroup%
\sethlcolor{slightgray}%
\hl{\ttfamily #1}%
\endgroup
}

%%%%%%%%%%%%%%%%%%%%%%%%%%%%%%%%%%%%%%%%%%%
%
% Kotlin syntax highlighting
%
%%%%%%%%%%%%%%%%%%%%%%%%%%%%%%%%%%%%%%%%%%%

\usepackage[skins,breakable,listings]{tcolorbox}

\usepackage[dvipsnames]{xcolor}
\usepackage[table]{xcolor}
\lstdefinelanguage{kotlin}{
comment=[l]{//},
commentstyle={\color{gray}\ttfamily},
emph={delegate, filter, firstOrNull, forEach, it, lazy, mapNotNull, println, @Repeat, return@},
emphstyle={\color{OrangeRed}},
identifierstyle=\color{black},
keywords={abstract, actual, as, as?, break, by, class, companion, continue, data, do, dynamic, else, enum, expect, false, final, for, fun, get, if, import, in, infix, interface, internal, is, null, object, open, operator, override, package, private, public, return, sealed, set, super, suspend, this, throw, true, try, typealias, val, var, vararg, when, where, while, tailrec, reified},
keywordstyle={\color{blue}\bfseries},
morecomment=[s]{/*}{*/},
morestring=[b]",
morestring=[s]{"""*}{*"""},
ndkeywords={@Deprecated, @JvmField, @JvmName, @JvmOverloads, @JvmStatic, @JvmSynthetic, Array, Byte, Double, Float, Int, Integer, Iterable, Long, Runnable, Short, String},
ndkeywordstyle={\color{BurntOrange}\bfseries},
sensitive=true,
stringstyle={\color{ForestGreen}\ttfamily},
literate={`}{{\char0}}1
}

%%%%%%%%%%%%%%%%%%%%%%%%%%%%%%%%%%%%%%%%%%%
%
% Color boxes
%
%%%%%%%%%%%%%%%%%%%%%%%%%%%%%%%%%%%%%%%%%%%

\tcbset{
enhanced jigsaw,
breakable,
listing only,
boxsep=-1pt,
top=-1pt,
bottom=-0.5pt,
right=-0.5pt,
overlay first={
\node[black!50] (S) at (frame.south) {\Large\ding{34}};
\draw[dashed,black!50] (frame.south west) -- (S) -- (frame.south east);
},
overlay middle={
\node[black!50] (S) at (frame.south) {\Large\ding{34}};
\draw[dashed,black!50] (frame.south west) -- (S) -- (frame.south east);
\node[black!50] (S) at (frame.north) {\Large\ding{34}};
\draw[dashed,black!50] (frame.north west) -- (S) -- (frame.north east);
},
overlay last={
\node[black!50] (S) at (frame.north) {\Large\ding{34}};
\draw[dashed,black!50] (frame.north west) -- (S) -- (frame.north east);
},
before={\par\vspace{10pt}},
after={\par\vspace{\parskip}\noindent}
}

\newtcblisting{kotlinlisting}[1][]{%
width=20.5cm,
left=20pt,
top=5pt,
listing options={
language=kotlin,
basicstyle=\ttfamily\normalsize,
%numberstyle=\footnotesize,
showstringspaces=false,
tabsize=2,
breaklines=true,
numbers=none,
inputencoding=utf8,
escapeinside={(*}{*)},
#1
},
underlay unbroken and first={%
\path[draw=none] (interior.north west) rectangle node[white]{\includegraphics[width=10mm]{../figures/kotlin_file.png}} ([xshift=-18mm,yshift=-20mm]interior.north west);
}
}

\newtcblisting{pythonlisting}[1][]{
width=17cm,
left=20pt,
top=5pt,
listing options={
language=Python,
basicstyle=\ttfamily\normalsize,
upquote=true,
breaklines=true,
showstringspaces=false,
keywordstyle=\color{blue}\bfseries,
escapeinside={(*}{*)},
#1
},
fonttitle=\ttfamily\small,
underlay unbroken and first={
\path[draw=none] (interior.north west) rectangle node[white]{\includegraphics[width=10mm]{../figures/python_icon.png}} ([xshift=-18mm,yshift=-20mm]interior.north west);
}
}

% Imitate syntax error
\usepackage{ulem}
\makeatletter
\def\uwave{\bgroup \markoverwith{\lower7.5\p@\hbox{\sixly \textcolor{red}{\char58}}}\ULon}
\font\sixly=lasy6 % does not re-load if already loaded, so no memory problem.
\makeatother

\usepackage{tikz}
\usepackage[skins,breakable,listings]{tcolorbox}
\usepackage{pgfplots}
\usepackage{tikz-qtree}
\usepackage{graphicx}

\usepackage{include/preamble}


% Custom notation
\newcommand{\fdeep}{\vf^{(1:L)}}
\newcommand{\flast}{\vf^{(L)}}
\newcommand{\Jx}{J_{\vx \rightarrow \vy}}
\newcommand{\Jxx}{J_{\vx \rightarrow \vy}(\vx)}
\newcommand{\Jy}{J_{\vy \rightarrow \vx}}
\newcommand{\Jyy}{J_{\vy \rightarrow \vx}(\vy)}
\newcommand{\detJyy}{ \left| J_{\vy \rightarrow \vx}(\vy) \right|}

\newcommand\transpose{{\textrm{\tiny{\sf{T}}}}}
\newcommand{\note}[1]{}
\newcommand{\hlinespace}{~\vspace*{-0.15cm}~\\\hline\\\vspace*{0.15cm}}
\newcommand{\embeddingletter}{g}
\newcommand{\bo}{{\sc bo}}
\newcommand{\agp}{Arc \gp}

\newcommand{\D}{\mathcal{D}}
\newcommand{\X}{\mathbf{X}}
\newcommand{\y}{y}
\newcommand{\data} {\X, \y}
\newcommand{\x}{\mathbf{x}}
\newcommand{\f}{\mathit{f}}

\newcommand{\fx}{ f(\mathbf{x}) }
\newcommand{\U}{\mathcal{U}}
\newcommand{\E}{\mathbf{E}}


\newcommand{\bardist}[0]{\hspace{-0.2cm}}

\newlength{\arrowsize}
\pgfarrowsdeclare{biggertip}{biggertip}{
\setlength{\arrowsize}{10pt}
\addtolength{\arrowsize}{2\pgflinewidth}
\pgfarrowsrightextend{0}
\pgfarrowsleftextend{-5\arrowsize}
}{
\setlength{\arrowsize}{1pt}
\addtolength{\arrowsize}{\pgflinewidth}
\pgfpathmoveto{\pgfpoint{-5\arrowsize}{4\arrowsize}}
\pgfpathlineto{\pgfpointorigin}
\pgfpathlineto{\pgfpoint{-5\arrowsize}{-4\arrowsize}}
\pgfusepathqstroke
}


% Custom commmands.

\def\jointspacing{\vspace{0.3in}}

\def\boxwidth{0.21\columnwidth}
\newcommand{\gpdrawbox}[1]{
\setlength\fboxsep{0pt}
\hspace{-0.36in}
\fbox{\hspace{-4mm}
%\includegraphics[width=\boxwidth]{../figures/deep_draws/deep_gp_sample_layer_#1}
\hspace{-4mm}}}

\newcommand{\mappic}[1]{
%\hspace{-0.05in}\includegraphics[width=\boxwidth]{../../figures/seed-0-map/latent_coord_map_layer_#1}
}

\newcommand{\mappiccon}[1]{
%\hspace{-0.05in}\includegraphics[width=\boxwidth]{../../figures/seed-0-map-connected/latent_coord_map_layer_#1}
}

\newcommand{\spectrumpic}[1]{
%\includegraphics[trim=4.5mm 0mm 4mm 3mm, clip, width=0.44\columnwidth]{../figures/spectrum/layer-#1}
}

\usepackage{dsfont}

\newcommand{\feat}{\vh}
\newcommand{\bs}{\blacksquare}
\newcommand{\ws}{\square}

\tikzset{
  treenode/.style = {shape=rectangle, rounded corners,
  draw, align=center,
  top color=white, bottom color=blue!20},
  root/.style     = {treenode, font=\tiny, bottom color=red!30},
  env/.style      = {treenode, font=\tiny},
  dummy/.style    = {circle,draw}
}


\begin{document}
  \begin{poster}
    %%% Header
    \begin{center}
      \begin{pcolumn}{1.03}
        %%% Title
        \begin{minipage}[c][9cm][c]{0.85\textwidth}
          \begin{center}
          {\veryHuge \textbf{A Tree Sampler for Bounded Context-Free Languages}}\\[10mm]
          {\huge Breandan Considine}
          \end{center}
        \end{minipage}
      \end{pcolumn}
    \end{center}

    \vspace*{-0.5cm}

    \large


    %%%%%%%%%%%%%%%%%%%%%%%%%%%%%%%%%%%%%%%%%%%%%%%%%%%%%%%%%%%%%%%%%%%%%%
    %%% Beginning of Document
    %%%%%%%%%%%%%%%%%%%%%%%%%%%%%%%%%%%%%%%%%%%%%%%%%%%%%%%%%%%%%%%%%%%%%%

    \Large

    \begin{multicols}{2}


      \mysection{Main Idea}

      \vspace*{-1cm}
      \null\hspace*{3cm}\begin{minipage}[c]{0.85\columnwidth}
      \begin{itemize}
%        \item Sampling with rejection is unnecessary if you can map onto a simpler distribution
        \item Analytic combinatorics: if you can count it, then you can sample from it!
        \item We implement a bijection between labeled binary trees in BCFLs and $\mathbb{Z}_{|T|}$
        \item Allows for communication-free parallel no-replacement sampling in $\widetilde{\mathcal{O}}(1)$
      \end{itemize}
      \end{minipage}

      \jointspacing

      \mysection{Semiring Parsing}
      \null\hspace*{3cm}\begin{minipage}[c]{0.85\columnwidth}
          Given a CFG $\mathcal{G} = \langle V, \Sigma, P, S\rangle$ in Chomsky Normal Form (CNF), we may construct a recognizer $R_\mathcal{G}: \Sigma^n \rightarrow \mathbb{B}$ for strings $\sigma: \Sigma^n$ as follows. Let $2^V$ be our domain, where $0$ is $\varnothing$, $\oplus$ is $\cup$, and $\otimes$ be defined as:\vspace{0.5cm}
      \end{minipage}

      \[
        s_1 \otimes s_2 = \{C \mid \langle A, B\rangle \in s_1 \times s_2, (C\rightarrow AB) \in P\}
      \]

      \null\hspace*{3cm}\begin{minipage}[c]{0.85\columnwidth}
If we define $\hat\sigma_r = \{w \mid (w \rightarrow \sigma_r) \in P\}$, then construct a matrix with unit nonterminals on the superdiagonal, $M_0[r+1=c](G', \sigma) = \;\hat\sigma_r$ the fixpoint $M_{i+1} = M_i + M_i^2$ is fully determined by the first diagonal:\vspace{0.5cm}
\end{minipage}

\begin{align*}
M_0=
\begin{pNiceMatrix}[nullify-dots,xdots/line-style=loosely dotted]
   \varnothing & \hat\sigma_1 & \varnothing & \Cdots & \varnothing \\
   \Vdots      & \Ddots   & \Ddots      & \Ddots & \Vdots\\
               &          &             &        & \varnothing\\
               &          &             &        & \hat\sigma_n \\
   \varnothing & \Cdots   &             &        & \varnothing
\end{pNiceMatrix} \Rightarrow
\begin{pNiceMatrix}[nullify-dots,xdots/line-style=loosely dotted]
  \varnothing & \hat\sigma_1 & \Lambda & \Cdots & \varnothing \\
  \Vdots      & \Ddots   & \Ddots  & \Ddots & \Vdots\\
              &          &         &        & \Lambda\\
              &          &         &        & \hat\sigma_n \\
  \varnothing & \Cdots   &         &        & \varnothing
\end{pNiceMatrix} \Rightarrow \ldots \Rightarrow M_\infty =
\begin{pNiceMatrix}[nullify-dots,xdots/line-style=loosely dotted]
   \varnothing & \hat\sigma_1 & \Lambda & \Cdots & \Lambda^*_\sigma\\
   \Vdots      & \Ddots   & \Ddots  & \Ddots & \Vdots\\
               &          &         &        & \Lambda\\
               &          &         &        & \hat\sigma_n \\
   \varnothing & \Cdots   &         &        & \varnothing
\end{pNiceMatrix}
\end{align*}

\null\hspace*{3cm}\begin{minipage}[c]{0.85\columnwidth}
CFL membership is recognized by $R(G', \sigma) = [S \in \Lambda^*_\sigma] \Leftrightarrow [\sigma \in \mathcal{L}(G)]$.
\end{minipage}

      \jointspacing

      \mysection{Parsing Dynamics}

      \null\hspace*{3cm}\begin{minipage}[c]{0.85\columnwidth}
      Let us consider an example with two holes, $\sigma = 1$ \underline{\hspace{1cm}} \underline{\hspace{1cm}}, and the grammar being $G=\{S\rightarrow N O N, O \rightarrow + \mid \times, N \rightarrow 0 \mid 1\}$. This can be rewritten into CNF as $G'=\{S \rightarrow N L, N \rightarrow 0 \mid 1, O \rightarrow \times \mid +, L \rightarrow O N\}$.\\
      \end{minipage}

      \null\hspace*{3cm}\begin{minipage}[c]{\columnwidth}
                      \resizebox{.85\columnwidth}{!}{
                      {\renewcommand{\arraystretch}{1.2}
\begin{tabular}{|c|c|c|c|}
  \hline
  & $2^V$ & $\mathbb{B}^{|V|}$ & $\mathbb{B}^{|V|}\rightarrow\mathbb{B}^{|V|}$\\\hline
  $M_0$ & \begin{pmatrix}
  \phantom{V} & \tiny{\{N\}} &         &             \\
              &              & \{N,O\} &             \\
              &              &         & \{N,O\} \\
              &              &         &
  \end{pmatrix} & \begin{pmatrix}
  \phantom{V} & \ws\bs\ws\ws &              &              \\
              &              & \ws\bs\bs\ws &              \\
              &              &              & \ws\bs\bs\ws \\
              &              &              &
  \end{pmatrix} & \begin{pmatrix}
     \phantom{V} & V_{0, 1} &          &          \\
                 &          & V_{1, 2} &          \\
                 &          &          & V_{2, 3} \\
                 &          &          &
  \end{pmatrix} \\\hline
  $M_1$ & \begin{pmatrix}
  \phantom{V} & \tiny{\{N\}} & \varnothing &         \\
              &              & \{N,O\}     & \{L\}   \\
              &              &             & \{N,O\} \\
              &              &             &
  \end{pmatrix} & \begin{pmatrix}
  \phantom{V} & \ws\bs\ws\ws & \ws\ws\ws\ws &              \\
              &              & \ws\bs\bs\ws & \bs\ws\ws\ws \\
              &              &              & \ws\bs\bs\ws \\
              &              &              &
  \end{pmatrix} & \begin{pmatrix}
  \phantom{V} & V_{0, 1} & V_{0, 2} &          \\
              &          & V_{1, 2} & V_{1, 3} \\
              &          &          & V_{2, 3} \\
              &          &          &
  \end{pmatrix} \\\hline
  $M_\infty$ & \begin{pmatrix}
  \phantom{V} & \tiny{\{N\}} & \varnothing & \{S\}   \\
              &              & \{N,O\}     & \{L\}   \\
              &              &             & \{N,O\} \\
              &              &             &
  \end{pmatrix} & \begin{pmatrix}
  \phantom{V} & \ws\bs\ws\ws & \ws\ws\ws\ws & \ws\ws\ws\bs \\
              &              & \ws\bs\bs\ws & \bs\ws\ws\ws \\
              &              &              & \ws\bs\bs\ws \\
              &              &              &
  \end{pmatrix} & \begin{pmatrix}
  \phantom{V} & V_{0, 1} & V_{0, 2} & V_{0, 3} \\
              &          & V_{1, 2} & V_{1, 3} \\
              &          &          & V_{2, 3} \\
              &          &          &
  \end{pmatrix}\\\hline
\end{tabular}\\
}
                      }\\
      \end{minipage}

      \null\hspace*{3cm}\begin{minipage}[c]{0.90\columnwidth}
      This procedure decides if $\exists \sigma' \in \mathcal{L}(G) \mid \sigma' \sqsubseteq \sigma$ but forgets provenance.
      \end{minipage}
      \jointspacing

      \mysection{Encoding CFL Sketching into SAT}
      \null\hspace*{3cm}\begin{minipage}[c]{0.90\columnwidth}
      \begin{itemize}
        \item CYK parser can be lowered onto a Boolean tensor $\mathbb{B}^{n\times n \times |V|}$ (Valiant, 1975)
        \item Binarized CYK parser can be compiled to SAT to solve for $\mathbf{M}^*$ directly
%        \item Enables sketch-based synthesis in either $\sigma$ or $\mathcal G$: just use variables for holes!
        \item We simply encode the characteristic function, i.e., $\mathds{1}_{\subseteq V}: 2^V\rightarrow \mathbb{B}^{|V|}$
        \item $\oplus, \otimes$ are defined as $\boxplus, \boxtimes$, so that the following diagram commutes:
        \[\begin{tikzcd}
            2^V \times 2^V \arrow[r, "\oplus/\otimes"] \arrow[d, "\mathds{1}^2"]
            & 2^V \arrow[d, "\mathds{1}\phantom{^{-1}}"] \\
            \mathbb{B}^{|V|} \times \mathbb{B}^{|V|} \arrow[r, "\boxplus/\boxtimes", labels=below] \arrow[u, "\mathds{1}^{-2}"]
            & \mathbb{B}^{|V|} \arrow[u, "\mathds{1}^{-1}"]
        \end{tikzcd}\]
        \item These operators can be lifted into matrices and tensors in the usual manner
%        \item In most cases, only a few nonterminals will be active at any given time
%        \item If density is desired, possible to use the Maculay representation
%        \item If you know of a more efficient encoding, please let us know!
      \end{itemize}
      \vspace{1.8cm}
      \end{minipage}

      \jointspacing

      \pagebreak      \mysection{A Nested Datatype for BCFLs}

      \jointspacing

      \hspace*{2cm}\begin{minipage}[c]{0.90\columnwidth}
      We define a map from nonterminals $\mathbb{T}_3 = (V \cup \Sigma) \rightharpoonup \mathbb{T}_2$ onto the datatype $\mathbb{T}_2 = (V \cup \Sigma) \times (\mathbb{N} \rightharpoonup \mathbb{T}_2\times\mathbb{T}_2)$, whose inhabitants satisfy the recurrence:
      \end{minipage}

      \begin{equation*}
        L(p) = 1 + p L(p) \phantom{addspace} P(a) = V + a L\big(V^2P(a)^2\big)
      \end{equation*}

      \hspace*{2cm}\begin{minipage}[c]{0.90\columnwidth}
        Each $\mathbb{T}_2$ consists of a root nonterminal, and a list of distinct products, e.g.,
      \end{minipage}

      \jointspacing

      \hspace*{10cm}\resizebox{0.5\columnwidth}{!}{
        \begin{tikzpicture}
        [
          grow                    = right,
          sibling distance        = 3em,
          level distance          = 5em,
          edge from parent/.style = {draw, -latex},
          every node/.style       = {font=\footnotesize},
          sloped
        ]
          \node [root] {S}
          child { node [env] {BC}
          child { node [root] {B}
          child { node [env] {RD}
          child { node [root] {R} edge from parent node [above] }
          child { node [root] {D} edge from parent node [above] }
          edge from parent node [above] }
          edge from parent node [below] }
          child { node [root] {C}
          child { node [env] {\ldots\vphantom{BB}} edge from parent node [below] }
%  child { edge from parent node [above] {\ldots} }
          edge from parent node [below] }
          edge from parent node [above] }
          child { node [env] {\ldots\vphantom{BB}} edge from parent node [below] }
          child { node [env] {AB}
          child { node [root] {A}
          child {
            node [env] {QC}
            child { node [root] {Q} edge from parent node [above] }
            child { node [root] {C} edge from parent node [above] }
            edge from parent node [above]
          }
%    child { node [env] {ZQ} edge from parent node [above] }
          child { node [env] {\ldots\vphantom{BB}} edge from parent node [below] }
          edge from parent node [below] }
          child { node [root] {B}
          child { node [env] {RD}
          child { node [root] {R} edge from parent node [above] }
          child { node [root] {D} edge from parent node [above] }
          edge from parent node [above] }
          edge from parent node [below] }
          edge from parent node [above] };
        \end{tikzpicture}
      }

      \vspace{1.5cm}

      \hspace*{2cm}\begin{minipage}[c]{0.90\columnwidth}
      Morally, $\mathbb{T}_2$ represents an implicit set of possible trees sharing the same root, where $\mathbb{T}_3$ is a dictionary of possible $\mathbb{T}_2$ values indexed by possible roots, given by a specific CFG under a porous string. Instead of $\hat\sigma_r$ we initialize using $\Lambda(\sigma_r)$:
      \end{minipage}

      \[
           \Lambda(s: \underline\Sigma) \mapsto \begin{cases}
               \bigoplus_{s\in \Sigma} \Lambda(s) & \text{if $s$ is a hole,} \vspace{5pt}\\
               \big\{\mathbb{T}_2\big(w, \big[\langle\mathbb{T}_2(s), \mathbb{T}_2(\varepsilon)\rangle\big]\big) \mid (w \rightarrow s)\in P\big\} & \text{otherwise.}
           \end{cases}
      \]
%
      \hspace*{2cm}\begin{minipage}[c]{0.90\columnwidth}
      The operations $\oplus, \otimes: \mathbb{T}_3 \times \mathbb{T}_3 \rightarrow \mathbb{T}_3$ are then redefined over trees as follows:
      \end{minipage}

      \[
          X \oplus Z \mapsto \bigcup_{\mathclap{k\in \pi_1(X \cup Z)}}\Big\{k \Rightarrow \mathbb{T}_2(k, x \cup z) \mid x \in \pi_2(X\circ k), z \in \pi_2(Z\circ k)\Big\}
      \]

      \[
        X \otimes Z \mapsto \bigoplus_{\mathclap{(w\rightarrow xz) \in P}}\Big\{\mathbb{T}_2\Big(w, \big[\langle X\circ x, Z\circ z\rangle\big]\Big) \mid x \in \pi_1(X), z \in \pi_1(Z)\Big\}
      \]
%
%      \hspace*{2cm}\begin{minipage}[c]{0.90\columnwidth}
%      These operators group subtrees by their root nonterminal, then aggregate their children. Each $\Lambda$ becomes a dictionary indexed by the root nonterminal, which can be sampled by obtaining $(\Lambda_\sigma^* \circ S): \mathbb{T}_2$, then recursively choosing twins.
%      \end{minipage}

      \vspace{-0.5cm}
      \jointspacing

      \mysection{Sampling with Replacement}

      \jointspacing

      \hspace*{2cm}\begin{minipage}[c]{0.90\columnwidth}
      Given a PCFG whose productions indexed by each nonterminal are decorated with a probability vector $\mathbf{p}$, we define a tree sampler $\Gamma: \mathbb{T}_2 \rightsquigarrow \mathbb{T}$ like so:
      \end{minipage}

      \begin{equation*}
        \Gamma(T) \mapsto \begin{cases}
                            \text{Multi} \big(\texttt{children}(T), \mathbf{p}\big) & \text{ if $T$ is a \texttt{root}} \\
                            \big\langle \Gamma\big(\pi_1(T)\big), \Gamma\big(\pi_2(T)\big) \big\rangle & \text{ if $T$ is a \texttt{child} }
        \end{cases}
      \end{equation*}


      \hspace*{2cm}\begin{minipage}[c]{0.90\columnwidth}
      This relates to the generating function for the ordinary Boltzmann sampler,
      \end{minipage}

      \begin{equation*}
        \Gamma C(x) \mapsto \begin{cases}
                              \text{Bern} \left(\frac{A(x)}{A(x) + B(x)}\right) \rightarrow \Gamma A(x) \mid \Gamma B(x) & \text{ if } \mathcal{C}=\mathcal{A}+\mathcal{B} \\
                              \big\langle \Gamma A(x), \Gamma B(x)\big\rangle & \text{ if } \mathcal{C}=\mathcal{A} \times \mathcal{B}
        \end{cases}
      \end{equation*}


      \hspace*{2cm}\begin{minipage}[c]{0.90\columnwidth}
      however unlike Duchon et al. (2004), our work does require rejection to ensure exact-size sampling, as all trees contained in $\mathbb{T}_2$ are necessarily the same width.
      \end{minipage}

      \jointspacing

      \mysection{Sampling without Replacement}

      \jointspacing

      \hspace*{2cm}\begin{minipage}[c]{0.90\columnwidth}
      To sample all trees in a given $T: \mathbb{T}_2$ uniformly without replacement, we define a modular pairing function $\varphi: \mathbb{T}_2 \rightarrow \mathbb{Z}_{|T|} \rightarrow \texttt{BTree}$ using the construction:
      \end{minipage}

      \begin{equation*}
        \varphi(T: \mathbb{T}_2, i: \mathbb{Z}_{|T|}) \mapsto \begin{cases}
          \Big\langle\texttt{BTree}\big(\texttt{root}(T)\big), i\Big\rangle & \text{if $T$ is a leaf,} \vspace{5pt}\\
          \text{Let } b = |\texttt{children}(T)|,\\
          \phantom{\text{Let }} q_1, r=\big\langle\lfloor\frac{i}{b}\rfloor, i \pmod{b}\big\rangle,\\
          \phantom{\text{Let }} lb, rb = \texttt{children}[r],\\
          \phantom{\text{Let }} T_1, q_2 = \varphi(lb, q_1),\\
          \phantom{\text{Let }} T_2, q_3 = \varphi(rb, q_2) \text{ in } \\
          \Big\langle\texttt{BTree}\big(\texttt{root}(T), T_1, T_2\big), q_3\Big\rangle & \text{otherwise.} \\
        \end{cases}
      \end{equation*}

      \hspace*{2cm}\begin{minipage}[c]{0.90\columnwidth}
      Then instead of sampling trees, we can simply sample integers WOR from $\mathbb{Z}_{|T|}$.\vspace{1.8cm}
      \end{minipage}

      \jointspacing

    \end{multicols}

    \vspace{-2.2cm}
    \bottombox{
    %% QR code
    %    \hfill\bottomboxlogo{img/kotlin_logo.png}
    % Comment out the line below out to hide logo

    \hspace{1.8cm}
    \begin{minipage}[c][0.1\paperheight][c]{0.18\textwidth}\qrcode[height=2.6in]{https://tidyparse.github.io/} \end{minipage}
    \begin{minipage}[c][0.1\paperheight][c]{0.25\textwidth}\includegraphics[height=2.6in]{../figures/tidyparse_logo.png} \end{minipage}
    \hspace{-4cm}
    \begin{minipage}[c][0.1\paperheight][c]{0.33\textwidth}\includegraphics[height=3in]{../figures/mcgill.png} \end{minipage}
    \hspace{2cm}
    \begin{minipage}[c][0.1\paperheight][c]{0.33\textwidth}\includegraphics[height=3.2in]{../figures/mila.png} \end{minipage}
    %    \hfill\bottomboxlogo{img/mila_mauve.png} % \hfill shifts the logo across so it meets the right hand side margin
    % Note that \bottomboxlogo takes an optional width argument. It defaults to the following:
    % \hfill\bottomboxlogo[width=\textwidth]{<path_to_image_file>}
    % where \textwidth is actually the width of a minipage which is defined in the \bottombox command of
    % betterportaitposter.cls It's a standard \includegraphics command in there, so easy to change if
    % you need to add a border etc.
    }
\end{poster}
\end{document}
